\section{无穷直积}
%@see: 《Elements of Set Theory》 P54 INFINITE CARTESIAN PRODUCTS
我们在前面学习了有限个集合的直积，
现在利用映射统一表示任意指标集上的直积.
取集合\(I\)作为指标集，
设\(H\)是映射且\(I\subseteq\dom H\).
定义
\begin{align*}
    \BigTimes_{i\in I}H(i)
    &\defeq\Set{f\given f\text{是映射}\land\dom f=I
    \land(\forall i\in I)[f(i)\in H(i)]}.
\end{align*}
其元素称为\(I\)元组（\(I\)-tuple），
第\(i\)个坐标就是\(f(i)\).
这里不要求\(I\)是无穷集.

\begin{proposition}
    对任意集合\(I\)和满足\(I\subseteq\dom H\)的映射\(H\)，
    上述直积是一个集合.
    \begin{proof}
        由子集公理在\(\ran H\)中构造
        \begin{align*}
            \mathscr{B}&\defeq\Set{Y\in\ran H\given(\exists i\in I)[Y=H(i)]},\\
            U&\defeq\bigcup\mathscr{B}.
        \end{align*}
        任意属于该直积的映射\(f\)都是\(I\times U\)的子集.
        因此直积可以写成
        \begin{align*}
            \BigTimes_{i\in I}H(i)
            &=\Set{f\in\Powerset(I\times U)\given
            f\text{是映射}\land\dom f=I
            \land(\forall i\in I)[f(i)\in H(i)]}.
        \end{align*}
        由直积存在定理、幂集公理和子集公理，
        右侧是集合.
        这个论证只证明直积的存在性，
        不声称它非空，
        因而不需要选择公理.
    \end{proof}
\end{proposition}

\begin{example}
    设\(A,I\)是集合，
    \(H=I\times\{A\}\)是以\(I\)为定义域的常值映射.
    证明：
    \begin{align*}
        \BigTimes_{i\in I}H(i)&=A^I.
    \end{align*}
    \begin{proof}
        对每个\(i\in I\)都有\(H(i)=A\).
        因而一个映射\(f\)属于左侧，
        当且仅当\(\dom f=I\)且对所有\(i\in I\)都有\(f(i)\in A\)，
        也就是\(f\colon I\to A\).
        这恰好是\(f\in A^I\)的定义，
        所以两集合相等.
    \end{proof}
\end{example}

\begin{property}
    若存在\(i_0\in I\)使\(H(i_0)=\emptyset\)，
    则\(\BigTimes_{i\in I}H(i)=\emptyset\).
    若\(I=\emptyset\)，
    则\(\BigTimes_{i\in I}H(i)=\{\emptyset\}\).
    \begin{proof}
        在第一种情形中，
        若\(f\)属于直积，
        就必须有\(f(i_0)\in\emptyset\)，
        矛盾，所以直积为空集.
        在第二种情形中，
        唯一以\(\emptyset\)为定义域的映射是空映射\(\emptyset\).
        条件\((\forall i\in\emptyset)[f(i)\in H(i)]\)成立，
        所以空映射属于直积，
        且没有其他元素.
        特别地，\(A^{\emptyset}=\{\emptyset\}\)，
        即使\(A=\emptyset\)也如此.
    \end{proof}
\end{property}

%@see: 《Elements of Set Theory》 P55
反过来，若每个\(H(i)\)都非空，
要得到直积中的一个元素，
就必须给出一个同时为每个指标选择一个元素的映射.
对于任意指标集的这项断言，
选择公理给出如下保证.
\begin{axiom}[选择公理(第二种形式)]
    %@see: 《Elements of Set Theory》 P55 Axiom of Choice (second form)
    对任意集合\(I\)及以\(I\)为定义域的映射\(H\)，
    如果\((\forall i\in I)[H(i)\neq\emptyset]\)，
    那么\(\BigTimes_{i\in I}H(i)\neq\emptyset\).
\end{axiom}

\section{选择公理}
\begin{definition}
    %@see: 《点集拓扑讲义（第四版）》（熊金城） P36 定义1.8.1
    设\(X\)是集合，
    \(\tilde{X}\defeq\Powerset X-\{\emptyset\}\).
    如果映射\(\epsilon\colon\tilde{X}\to X\)满足
    \begin{align*}
        (\forall A\in\tilde{X})[\epsilon(A)\in A],
    \end{align*}
    则称“\(\epsilon\)是集合\(X\)的一个\DefineConcept{选择函数}”.
\end{definition}

\begin{axiom}[选择公理(第三种形式)]
    %@see: 《点集拓扑讲义（第四版）》（熊金城） P36 公理1.8.1
    %@see: 《Elements of Set Theory》 P151 Theorem 6M (3)
    任何集合都有选择函数.
\end{axiom}

\begin{theorem}\label{theorem:选择公理.等价形式1}
    %@see: 《点集拓扑讲义（第四版）》（熊金城） P37 定理1.8.2
    设\(\mathscr{A}\)是由非空集合构成的集族，
    则存在映射\(\nu\colon\mathscr{A}\to\bigcup\mathscr{A}\)，
    使得
    \begin{align*}
        (\forall A\in\mathscr{A})[\nu(A)\in A].
    \end{align*}
    \begin{proof}
        令\(X=\bigcup\mathscr{A}\).
        由选择公理的第三种形式，
        取\(X\)的选择函数\(\epsilon\).
        每个\(A\in\mathscr{A}\)都是\(X\)的非空子集，
        因而\(\mathscr{A}\subseteq\dom\epsilon\).
        令\(\nu\)为\(\epsilon\)在\(\mathscr{A}\)上的限制，
        则\(\nu(A)=\epsilon(A)\in A\).
        当\(\mathscr{A}=\emptyset\)时，
        所得\(\nu\)为空映射，
        同样满足条件.
    \end{proof}
\end{theorem}

\begin{proposition}
    \cref{theorem:选择公理.等价形式1}中的断言、选择公理的第二种形式和第三种形式互相等价.
    \begin{proof}
        第三种形式推出集族形式的证明已在上述定理中给出.

        下面由集族形式推出第二种形式.
        设\(H\)的定义域为\(I\)，
        且每个\(H(i)\)都非空.
        对集族\(\mathscr{A}=\Set{H(i)\given i\in I}\)应用集族形式，
        得到\(\nu\)满足\(\nu(H(i))\in H(i)\).
        定义映射
        \begin{align*}
            f\colon I\to\bigcup\mathscr{A},\quad
            i&\mapsto\nu(H(i)).
        \end{align*}
        它的图像是\(I\times\bigcup\mathscr{A}\)的子集，
        并且\(f(i)\in H(i)\)，
        所以\(f\in\BigTimes_{i\in I}H(i)\).
        即使不同指标对应同一个集合，
        这个构造仍然成立，
        因为各指标可以选取同一个元素.

        最后由第二种形式推出第三种形式.
        给定集合\(X\)，
        令\(I=\Powerset X-\{\emptyset\}\)，
        并令\(H\)为\(I\)上的恒等映射，
        即\(H(B)=B\).
        每个\(H(B)\)都非空，
        所以存在\(\epsilon\in\BigTimes_{B\in I}H(B)\).
        于是\(\dom\epsilon=I\)，
        且\(\epsilon(B)\in B\subseteq X\)，
        因此\(\epsilon\)就是\(X\)的选择函数.
        当\(X=\emptyset\)时，
        指标集为空，
        其直积含有唯一的空映射，
        因而结论仍成立.
    \end{proof}
\end{proposition}

\begin{definition}
    %@see: 《点集拓扑讲义（第四版）》（熊金城） P37 定义1.8.2
    设\(\mathscr{F}\)是集族.
    如果
    \begin{align*}
        (\forall A,B\in\mathscr{F})[A\subseteq B\lor B\subseteq A],
    \end{align*}
    则称“\(\mathscr{F}\)是一个\DefineConcept{套}（chain）”.
    空集族也满足这个全称条件，
    因而是一个套.
\end{definition}

\begin{definition}
    设\(M\in\mathscr{F}\).
    若
    \begin{align*}
        (\forall B\in\mathscr{F})[M\subseteq B\implies B=M],
    \end{align*}
    则称\(M\)是\(\mathscr{F}\)关于\(\subseteq\)的\DefineConcept{极大元}.
    若
    \begin{align*}
        (\forall B\in\mathscr{F})[B\subseteq M],
    \end{align*}
    则称\(M\)是\(\mathscr{F}\)关于\(\subseteq\)的\DefineConcept{最大元}.
    极大元只排除自身的真扩张，
    并不要求它与其他所有成员可比较.
\end{definition}

为避免在佐恩引理与图基引理之间循环论证，
先直接由选择函数证明一个辅助结论.
\begin{lemma}\label{lemma:无穷直积.单点扩张极大元}
    设非空集族\(\mathscr{C}\)满足：
    每个成员的任意子集仍属于\(\mathscr{C}\)，
    且\(\mathscr{C}\)中每个套的并仍属于\(\mathscr{C}\).
    则\(\mathscr{C}\)有关于\(\subseteq\)的极大元.
    \begin{proof}
        令\(U=\bigcup\mathscr{C}\)，
        由选择公理取\(U\)的选择函数\(c\).
        对每个\(B\in\mathscr{C}\)，令
        \begin{align*}
            E(B)&\defeq\Set{x\in U-B\given B\cup\{x\}\in\mathscr{C}}.
        \end{align*}
        当\(E(B)=\emptyset\)时令\(g(B)=B\)，
        否则令\(g(B)=B\cup\{c(E(B))\}\).
        于是\(g\colon\mathscr{C}\to\mathscr{C}\)，
        \(B\subseteq g(B)\)，
        且\(g(B)-B\)至多含一个元素.
        由非空性和对子集的封闭性可知\(\emptyset\in\mathscr{C}\).

        考虑满足以下三个条件的子族\(\mathscr{T}\subseteq\mathscr{C}\)：
        \begin{enumerate}
            \item \(\emptyset\in\mathscr{T}\)；
            \item 若\(B\in\mathscr{T}\)，则\(g(B)\in\mathscr{T}\)；
            \item 若\(\mathscr{D}\subseteq\mathscr{T}\)是套，
            则\(\bigcup\mathscr{D}\in\mathscr{T}\).
        \end{enumerate}
        所有这样的子族构成\(\Powerset\mathscr{C}\)的一个非空子集，
        因为\(\mathscr{C}\)本身满足三个条件.
        令\(\mathscr{T}_0\)为这些子族的交.
        每个参与求交的子族都含有\(\emptyset\)，
        所以\(\emptyset\in\mathscr{T}_0\).
        若\(B\in\mathscr{T}_0\)，
        则\(B\)属于每个这样的子族，
        故\(g(B)\)也属于每个这样的子族，
        从而\(g(B)\in\mathscr{T}_0\).
        若套\(\mathscr{D}\)包含于\(\mathscr{T}_0\)，
        它就包含于每个参与求交的子族，
        故\(\bigcup\mathscr{D}\)属于每个这样的子族，
        从而\(\bigcup\mathscr{D}\in\mathscr{T}_0\).
        因此\(\mathscr{T}_0\)仍满足三个条件，
        并且包含于任何满足这三个条件的子族.

        下面证明\(\mathscr{T}_0\)是套.
        令
        \begin{align*}
            \mathscr{D}_0&\defeq\Set{C\in\mathscr{T}_0\given
            (\forall B\in\mathscr{T}_0)[B\subseteq C\lor C\subseteq B]}.
        \end{align*}
        显然\(\emptyset\in\mathscr{D}_0\).
        固定\(C\in\mathscr{D}_0\)，令
        \begin{align*}
            \mathscr{T}_C&\defeq\Set{B\in\mathscr{T}_0\given
            B\subseteq C\lor g(C)\subseteq B}.
        \end{align*}
        我们验证\(\mathscr{T}_C\)满足上述三个条件.
        首先\(\emptyset\in\mathscr{T}_C\).
        对任意\(B\in\mathscr{T}_C\)，
        若\(g(C)\subseteq B\)，
        则\(g(C)\subseteq B\subseteq g(B)\)，
        故\(g(B)\in\mathscr{T}_C\).
        若\(B=C\)，结论也直接成立.
        剩下的情形是\(B\subsetneq C\).
        由于\(g(B)\in\mathscr{T}_0\)且\(C\in\mathscr{D}_0\)，
        \(g(B)\)与\(C\)可比较.
        若\(C\subseteq g(B)\)，
        由\(B\subsetneq C\subseteq g(B)\)和\(g(B)-B\)至多含一个元素，
        必有\(C=g(B)\).
        因此无论哪一种比较关系成立，
        都有\(g(B)\subseteq C\)，
        所以\(g(B)\in\mathscr{T}_C\).
        最后，若\(\mathscr{E}\subseteq\mathscr{T}_C\)是套，
        要么其某个成员包含\(g(C)\)，
        从而其并也包含\(g(C)\)；
        要么其所有成员都包含于\(C\)，
        从而其并包含于\(C\).
        又因\(\mathscr{T}_0\)对套的并封闭，
        可知\(\bigcup\mathscr{E}\in\mathscr{T}_C\).

        根据\(\mathscr{T}_0\)的最小性，
        有\(\mathscr{T}_0\subseteq\mathscr{T}_C\)，
        从而\(\mathscr{T}_C=\mathscr{T}_0\).
        因此每个\(B\in\mathscr{T}_0\)或者包含于\(C\subseteq g(C)\)，
        或者包含\(g(C)\)，
        这说明\(g(C)\in\mathscr{D}_0\).
        故\(\mathscr{D}_0\)对\(g\)封闭.

        若\(\mathscr{E}\subseteq\mathscr{D}_0\)是套，
        则\(V=\bigcup\mathscr{E}\in\mathscr{T}_0\).
        对任意\(B\in\mathscr{T}_0\)，
        若存在\(C\in\mathscr{E}\)满足\(B\subseteq C\)，
        则\(B\subseteq V\).
        否则由每个\(C\in\mathscr{E}\)都与\(B\)可比较，
        得到\(C\subseteq B\)，
        从而\(V\subseteq B\).
        因此\(V\in\mathscr{D}_0\)，
        即\(\mathscr{D}_0\)也对套的并封闭.
        由最小性再得\(\mathscr{T}_0\subseteq\mathscr{D}_0\)，
        故\(\mathscr{D}_0=\mathscr{T}_0\)，
        所以\(\mathscr{T}_0\)是套.

        令\(M=\bigcup\mathscr{T}_0\).
        由封闭性，\(M\in\mathscr{T}_0\)，
        进而\(g(M)\in\mathscr{T}_0\)，
        所以\(g(M)\subseteq M\).
        结合\(M\subseteq g(M)\)可得\(g(M)=M\)，
        因而\(E(M)=\emptyset\).
        若存在\(B\in\mathscr{C}\)使\(M\subsetneq B\)，
        取\(x\in B-M\).
        因为\(M\cup\{x\}\subseteq B\)，
        对子集的封闭性给出\(M\cup\{x\}\in\mathscr{C}\)，
        于是\(x\in E(M)\)，矛盾.
        所以\(M\)是\(\mathscr{C}\)的极大元.
    \end{proof}
\end{lemma}

\begin{theorem}[佐恩引理]\label{theorem:无穷直积.佐恩引理}
    %@see: 《Elements of Set Theory》 P151 Theorem 6M (6) Zorn's lemma
    设集族\(\mathscr{A}\)满足：
    每个包含于\(\mathscr{A}\)的套的并仍属于\(\mathscr{A}\)，
    即
    \begin{align*}
        (\forall\mathscr{B}\subseteq\mathscr{A})
        [\mathscr{B}\text{是套}\implies\bigcup\mathscr{B}\in\mathscr{A}],
    \end{align*}
    则\(\mathscr{A}\)中有关于\(\subseteq\)的极大元.
    \begin{proof}
        由于空集族是套，
        假设已经保证\(\emptyset=\bigcup\emptyset\in\mathscr{A}\)，
        特别地\(\mathscr{A}\neq\emptyset\).
        令
        \begin{align*}
            \mathscr{C}&\defeq\Set{\mathscr{B}\in\Powerset\mathscr{A}\given
            \mathscr{B}\text{是套}}.
        \end{align*}
        \(\mathscr{C}\)含有空集族，
        并且它的每个成员的任意子集仍是套，
        因而仍属于\(\mathscr{C}\).
        若\(\mathscr{E}\subseteq\mathscr{C}\)本身关于包含关系构成套，
        任取\(B_1,B_2\in\bigcup\mathscr{E}\)，
        存在\(\mathscr{B}_1,\mathscr{B}_2\in\mathscr{E}\)分别包含\(B_1,B_2\).
        两个子族中一个包含另一个，
        所以\(B_1,B_2\)同属于其中较大的一个套，
        从而\(B_1\subseteq B_2\)或\(B_2\subseteq B_1\).
        因此\(\bigcup\mathscr{E}\)也是套，
        即\(\bigcup\mathscr{E}\in\mathscr{C}\).

        对\(\mathscr{C}\)应用\cref{lemma:无穷直积.单点扩张极大元}，
        得到不能再扩大的套\(\mathscr{M}\subseteq\mathscr{A}\).
        由题设，\(M=\bigcup\mathscr{M}\in\mathscr{A}\).
        若\(B\in\mathscr{A}\)且\(M\subseteq B\)，
        则对每个\(C\in\mathscr{M}\)都有\(C\subseteq M\subseteq B\)，
        所以\(\mathscr{M}\cup\{B\}\)仍是\(\mathscr{A}\)中的套.
        根据\(\mathscr{M}\)的极大性，
        必有\(B\in\mathscr{M}\)，
        从而\(B\subseteq\bigcup\mathscr{M}=M\).
        于是\(B=M\)，
        这就证明\(M\)是\(\mathscr{A}\)的极大元.
    \end{proof}
\end{theorem}

\begin{definition}
    %@see: 《点集拓扑讲义（第四版）》（熊金城） P37 定义1.8.3
    设\(\mathscr{F}\)是集族.
    如果对于每个集合\(F\)都有
    \begin{align*}
        F\in\mathscr{F}
        &\iff F\text{的每个有限子集均属于}\mathscr{F},
    \end{align*}
    则称“\(\mathscr{F}\)是\DefineConcept{具有有限特征的}集族”.
    有限子集包括空集.
\end{definition}

\begin{example}
    佐恩引理的结论能否把“极大元”加强为“最大元”？
    \begin{solution}
        不能.
        取不同的元素\(a,b\)，令
        \begin{align*}
            \mathscr{A}&=\{\emptyset,\{a\},\{b\}\}.
        \end{align*}
        因为\(\{a\}\)与\(\{b\}\)不可比较，
        \(\mathscr{A}\)中的套不能同时含有这两个成员.
        因此每个套的并只能是\(\emptyset\)、\(\{a\}\)或\(\{b\}\)，
        均属于\(\mathscr{A}\)，
        所以它满足上述佐恩引理的假设.
        \(\{a\}\)和\(\{b\}\)都是极大元，
        但若有最大元，
        它必须同时包含\(a,b\)，
        而\(\mathscr{A}\)中没有这样的成员.
        这个集族还具有有限特征.
        若\(F\in\mathscr{A}\)，
        则\(F\)的每个子集均属于\(\mathscr{A}\)，
        特别地每个有限子集都属于\(\mathscr{A}\).
        反之，若\(F\)的每个有限子集属于\(\mathscr{A}\)，
        考察单元素子集可得\(F\subseteq\{a,b\}\).
        \(F\)不能同时含有\(a,b\)，
        否则有限子集\(\{a,b\}\)不属于\(\mathscr{A}\)，
        矛盾.
        所以\(F\)只能是\(\emptyset\)、\(\{a\}\)或\(\{b\}\)，
        因而\(F\in\mathscr{A}\).
        因此有限特征也不能保证最大元存在.
        即使额外要求所求成员包含指定的\(\emptyset\)，
        这个反例仍然适用.
    \end{solution}
\end{example}

\begin{lemma}\label{lemma:无穷直积.有限特征封闭性}
    %@see: 《点集拓扑讲义（第四版）》（熊金城） P37 引理1.8.3
    如果\(\mathscr{F}\)具有有限特征，
    则有：
    \begin{enumerate}
        \item \(\mathscr{F}\)的每个成员的任意子集都属于\(\mathscr{F}\)；
        \item 若\(\mathscr{F}\neq\emptyset\)，
        则\(\mathscr{F}\)中任意一个套的并都属于\(\mathscr{F}\).
    \end{enumerate}
    \begin{proof}
        对第一条，设\(B\subseteq A\in\mathscr{F}\).
        \(B\)的任意有限子集也是\(A\)的有限子集，
        所以属于\(\mathscr{F}\).
        再由有限特征的定义得到\(B\in\mathscr{F}\).

        对第二条，由非空性取\(A\in\mathscr{F}\)，
        再由第一条得到\(\emptyset\in\mathscr{F}\).
        设\(\mathscr{B}\subseteq\mathscr{F}\)是套，
        记\(V=\bigcup\mathscr{B}\).
        任取\(V\)的有限子集\(K\).
        若\(K=\emptyset\)，则\(K\in\mathscr{F}\).
        若\(K\neq\emptyset\)，
        将其元素列为\(x_1,\dotsc,x_n\)，
        分别取\(B_1,\dotsc,B_n\in\mathscr{B}\)使\(x_j\in B_j\).
        这只是有限次选取，
        可逐次完成.
        在套的这有限个成员中逐次比较包含关系，
        得到一个包含其余所有成员的\(B_k\).
        于是\(K\subseteq B_k\in\mathscr{F}\)，
        由第一条可得\(K\in\mathscr{F}\).
        因为\(V\)的每个有限子集都属于\(\mathscr{F}\)，
        有限特征给出\(V\in\mathscr{F}\).
        当\(\mathscr{B}=\emptyset\)时，
        \(V=\emptyset\in\mathscr{F}\)，
        因而空套也已包含在证明中.
    \end{proof}
\end{lemma}

\begin{example}
    在上条引理的第二项中，
    能否去掉\(\mathscr{F}\neq\emptyset\)的条件？
    \begin{solution}
        不能.
        空集族\(\mathscr{F}=\emptyset\)具有有限特征：
        对任意集合\(F\)，
        \(F\in\mathscr{F}\)为假，
        而“\(F\)的每个有限子集属于\(\mathscr{F}\)”也为假，
        因为\(\emptyset\)是\(F\)的有限子集且\(\emptyset\notin\mathscr{F}\).
        因此定义中的等价式对每个\(F\)成立.
        但是\(\mathscr{B}=\emptyset\)是包含于\(\mathscr{F}\)的套，
        其并为\(\emptyset\)，
        而\(\emptyset\notin\mathscr{F}\).
        所以第二项若不加非空条件便不成立.
    \end{solution}
\end{example}

\begin{theorem}[图基引理]
    %@see: 《点集拓扑讲义（第四版）》（熊金城） P38 定理1.8.4
    非空的具有有限特征的集族中有关于\(\subseteq\)的极大元.
    \begin{proof}
        设\(\mathscr{F}\neq\emptyset\)具有有限特征.
        由\cref{lemma:无穷直积.有限特征封闭性}，
        \(\mathscr{F}\)中每个套的并仍属于\(\mathscr{F}\)，
        包括空套的并.
        因此\(\mathscr{F}\)满足\cref{theorem:无穷直积.佐恩引理}的假设，
        故有关于\(\subseteq\)的极大元.
    \end{proof}
\end{theorem}

\begin{corollary}
    %@see: 《点集拓扑讲义（第四版）》（熊金城） P40 推论1.8.5
    若\(\mathscr{F}\)具有有限特征且\(A\in\mathscr{F}\)，
    则\(\mathscr{F}\)有一个包含\(A\)的关于\(\subseteq\)的极大元.
    \begin{proof}
        令\(U=\bigcup\mathscr{F}\)，并构造
        \begin{align*}
            \mathscr{G}&\defeq\Set{C\in\Powerset(U-A)\given A\cup C\in\mathscr{F}}.
        \end{align*}
        由于\(A\in\mathscr{F}\)，
        有\(\emptyset\in\mathscr{G}\).
        设\(\mathscr{B}\subseteq\mathscr{G}\)是套.
        若\(\mathscr{B}=\emptyset\)，
        则\(\bigcup\mathscr{B}=\emptyset\in\mathscr{G}\).
        若\(\mathscr{B}\neq\emptyset\)，
        则\(\Set{A\cup C\given C\in\mathscr{B}}\)是\(\mathscr{F}\)中的套，
        由\cref{lemma:无穷直积.有限特征封闭性}可得
        \begin{align*}
            A\cup\bigcup\mathscr{B}
            &=\bigcup\Set{A\cup C\given C\in\mathscr{B}}\in\mathscr{F}.
        \end{align*}
        同时\(\bigcup\mathscr{B}\subseteq U-A\)，
        所以\(\bigcup\mathscr{B}\in\mathscr{G}\).
        由\cref{theorem:无穷直积.佐恩引理}，
        \(\mathscr{G}\)有极大元\(C_0\).
        令\(M=A\cup C_0\)，
        则\(A\subseteq M\in\mathscr{F}\).
        若\(M\subseteq B\in\mathscr{F}\)，
        则\(B-A\in\mathscr{G}\)且\(C_0\subseteq B-A\).
        极大性给出\(B-A=C_0\)，
        因而\(B=A\cup(B-A)=A\cup C_0=M\).
        所以\(M\)就是所求的极大元.
    \end{proof}
\end{corollary}
